% Part: propositional-logic

\documentclass[../../include/open-logic-part]{subfiles}

\begin{document}

\olpart{pl}{Önerme Mantığı}

\begin{editorial}
  Bu kısım klasik önerme mantığına ilişkin materyal içerir. İlk bölüm
  görece başlangıç düzeyindedir ve yalnızca tanımları ve sonuçları sıralar;
  birçok kanıt verilmez, alıştırma olarak bırakılır. Kanıt sistemleri ve
  tamlık teoremi hakkındaki materyal, ``FOL'' etiketi yanlış değerine
  ayarlanarak birinci mertebeden mantık kısmından içeri alınmıştır. Böylece
  yüklemler, terimler ve niceleyicilerle ilgili her şey dışarıda bırakılır ve
  $\Struct{M}$ ile gösterilen !!{structure}s hakkındaki anlatımın yerini
  $\pAssign{v}$ ile gösterilen !!{valuation}s hakkındaki anlatım alır.

  Bu kısmın daha ayrıntılı olacak biçimde genişletilmesi ve doğruluk
  fonksiyonları bakımından tamlık gibi başka konu ve sonuçların eklenmesi
  planlanmaktadır.
\end{editorial}

\olimport[syntax-and-semantics]{syntax-and-semantics}

\tagfalse{FOL}

\olimport[../first-order-logic/proof-systems]{proof-systems}

\iftag{prfSC}{%
  \olimport[../first-order-logic/sequent-calculus]{sequent-calculus}
}{}

\iftag{prfND}{%
  \olimport[../first-order-logic/natural-deduction]{natural-deduction}
}{}

\iftag{prfTab}{%
  \olimport[../first-order-logic/tableaux]{tableaux}
}{}

\iftag{prfAX}{%
  \olimport[../first-order-logic/axiomatic-deduction]{axiomatic-deduction}
}{}

\olimport[../first-order-logic/completeness]{completeness}

\tagtrue{FOL}

\OLEndPartHook

\end{document}
